\clearpage\fancyhead[L]{\small GAUSS AN VON LINDENAU.}\SourceLeaf{330}
\begin{center}\textbf{GAUSS AN VON LINDENAU.}\par\vspace{4pt}\hrule\vspace{4pt}\fontsize{10}{12}\selectfont Monatliche Correspondenz. Band XXIII. Seite 97. 98. 1811. Januar.\par\vspace{4pt}\hrule\vspace{4pt}\end{center}

\begin{flushright}Göttingen, den 14. Januar 1811.\end{flushright}
\par\indent --- --- --- An der hiesigen Sternwarte ließ sich in der Nacht vom 4. auf den 5. Januar erstmals seit beinahe zwei Monaten wieder versuchen, \emph{Pallas} zu finden. Da ich dazu nur bewegliche Instrumente verwenden konnte, erschwerten der Mondschein und die strenge Kälte (11$^\circ$ Réaum.) die Suche ungemein. Meine Bemühungen waren dennoch von Erfolg gekrönt. Ich fand \emph{Pallas} an der Stelle, an der sie nach meiner korrigierten Ephemeride stehen musste. Außerdem gelangen mir ein paar Vergleiche mit Sternen der \emph{Histoire céleste}. Daraus leitete ich unter dem Vorbehalt einer späteren genaueren Reduktion die folgende Position des Planeten ab.

\begin{center}\fontsize{10}{12}\selectfont\setlength{\tabcolsep}{4pt}\renewcommand{\arraystretch}{1.15}\begin{tabular}{c|c|c|c}
1811 & Mittlere Zeit & \shortstack{Rektaszension\\der Pallas} & \shortstack{Südliche\\Deklination}\\\hline
4. Januar & $15^{\text{u}}\;4^{\text{m}}\;42^{\text{s}}$ & $151^\circ\;40'\;5''$ & $22^\circ\;22'\;36''$\\
\end{tabular}\end{center}

\Needspace{16\baselineskip}\par\indent Meine erneut durchgesehenen Störungsrechnungen haben nun die folgenden Elemente der \emph{Pallas} ergeben. Sie weichen nur geringfügig von denen ab, die ich Ihnen zuletzt mitgeteilt habe (Dez. Heft 1810, S. 592) [S. 321 d. B.].

\begin{center}\fontsize{10}{12}\selectfont\setlength{\tabcolsep}{4pt}\renewcommand{\arraystretch}{1.15}\begin{tabular}{llr}
\shortstack[l]{Epoche der mittl. Länge\\im Meridian von Göttingen} & .\;.\;. & $126^\circ\;33'\;6''$\\
Länge des Perihels & .\;.\;.\;.\;.\;.\;.\;.\;.\;.\;.\;.\;.\;.\;. & $120\;55\;8$\\
Länge des aufsteigenden Knotens & .\;.\;.\;.\;.\;.\;.\;.\;.\;.\;. & $172\;32\;40$\\
Neigung der Bahn & .\;.\;.\;.\;.\;.\;.\;.\;.\;.\;.\;.\;.\;.\;.\;. & $34\;35\;14$\\
Tägliche mittlere tropische Bewegung & .\;.\;.\;.\;.\;.\;.\;.\;. & $768.980$\\
Exzentrizität & .\;.\;.\;.\;.\;.\;.\;.\;.\;.\;.\;.\;.\;.\;.\;.\;.\;. & $0.24160$\\
Logarithmus der großen Halbachse & .\;.\;.\;.\;.\;.\;.\;.\;. & $0.44278$\\
\end{tabular}\end{center}

\par\indent Hier folgt noch die berechnete Fortsetzung der Ephemeride [vom 3. Febr. bis 12. April 1811].
